\documentclass{article}

\usepackage[T1]{fontenc}

\usepackage{booktabs,graphicx,siunitx}

\begin{document}

\section{Mann-Whitney U}

\begin{tabular}{llllll}
\toprule
Group & N & Mean rank & Median & IQR & Sum of ranks \\
\midrule
Female & 120 & 124.30 & 5.00 & 3.00 & 14916.00 \\
Male & 120 & 116.70 & 4.00 & 4.00 & 14004.00 \\
\bottomrule
\end{tabular}

\begin{tabular}{llll}
\toprule
U & Z & p & r [95\% CI] \\
\midrule
7656 & 0.85 & .394 & 0.06 [$-$0.08, 0.21] \\
Hodges-Lehmann median difference = 0.00, 95\% CI [$-$0.00, 1.00] &  &  &  \\
\bottomrule
\end{tabular}

\includegraphics{figures/01_mann-whitney-u/figure_boxplot.png}

The nonparametric counterpart to the independent t-test, comparing whole distributions via ranks rather than means. A p below alpha means one group's values tend to be systematically higher or lower (stochastic dominance) - this equals a difference in medians only when the two groups have similar distribution shapes. r gives the effect size. Your significance threshold ($\alpha$) is 0.05. Method: Wilcoxon rank sum test with continuity correction.

A Mann-Whitney U test gave U=7656, Z=0.85, p = .394, r=.06 [$-$.08, .21].

Statistical basis: Mann-Whitney U test (Mann, H. B., Whitney, D. R. (1947). "On a test of whether one of two random variables is stochastically larger than the other." Annals of Mathematical Statistics, 18(1), 50-60.); Rank-biserial effect size r (Ben-Shachar MS, Lüdecke D, Makowski D (2020). "effectsize: Estimation of Effect Size Indices and Standardized Parameters." Journal of Open Source Software, 5(56), 2815.).

\end{document}